\lecture{8}{18. September 2025}{Structural Analysis -- Indeterminate trusses}

\begin{problemwithimage}{f5_19}{5-19}
  Determine the force developed in members $FE$, $EB$, and $BC$ of the truss and state if these members are in tension or compression.
\end{problemwithimage}

\begin{derivation}
  As we have no horizontal loads, we know that $A_x = 0$.
  Taking a moment equilibrium about $A$ yields:
  \begin{equation*}
    - \qty{11}{\kN} \cdot \qty{2}{\m} - \qty{22}{\kN} \cdot \qty{3.5}{\m} + D_y \cdot \qty{5,5}{\m} = 0 \implies D_y = \frac{\qty{11}{\kN} \cdot \qty{2}{\m} + \qty{22}{\kN} \cdot \qty{3.5}{\m} }{\qty{5.5}{\m} } = \qty{18}{\kN}
  .\end{equation*}
  A force equilibrium on the entire truss in the $y$ direction yields
  \begin{equation*}
    A_y + D_y - \qty{11}{\kN} - \qty{22}{\kN} = 0 \implies A_y = \qty{33}{\kN} - D_y = \qty{15}{\kN} 
  .\end{equation*}
  Hence, the support reactions are
  \begin{equation*}
    A_y = \qty{15}{\kN} \quad \text{and} \quad D_y = \qty{18}{\kN} 
  .\end{equation*}
  Now, we can start at joint $D$ and write a force equilibrium in the $y$ direction to get:
  \begin{equation*}
    \qty{18}{\kN} + F_{DE} \cdot \sin \ang{45} = 0 \implies F_{DE} = - \frac{\qty{18}{\kN}}{\sin \ang{45} }  = - \qty{25,46}{\kN}
  .\end{equation*}
  Which gives the equilibrium in the $x$ direction:
  \begin{equation*}
    - F_{CD} - F_{DE} \cdot \cos \ang{45} = 0 \implies F_{CD} = \qty{18}{\kN} 
  .\end{equation*}
  
  Moving to joint $C$, in the $x$ direction, we get
  \begin{equation*}
    F_{CD} - F_{BC} = 0 \implies F_{BC} = F_{CD} = \qty{18}{\kN} 
  .\end{equation*}
  And in the $y$ direction, we get
  \begin{equation*}
    F_{CE} - \qty{22}{\kN} = 0 \implies F_{CE} = \qty{22}{\kN} 
  .\end{equation*}
  
  For joint $A$, in the $y$ direction we get
  \begin{equation*}
    \qty{15}{\kN} + F_{AF} \sin \ang{45} = 0 \implies F_{AF} = - \frac{\qty{15}{\kN} }{\sin \ang{45}} = - \qty{21,21}{\kN} 
  .\end{equation*}
  In the $x$ direction at $A$ we get
  \begin{equation*}
    F_{AB} + F_{AF} \cdot \cos \ang{45} = 0 \implies F_{AB} = \qty{15}{\kN} 
  .\end{equation*}
  
  Now, we can move to joint $B$ in the $x$ direction and get
  \begin{align*}
    F_{BC} - F_{AB} + F_{BE} \cdot \num{0,6} &= 0 \\
    F_{BE} &= \frac{\qty{15}{\kN} - \qty{18}{\kN}}{\num{0.6} } = - \qty{5}{\kN}
  .\end{align*}

  Now, we move on to joint $F$ and use a horizontal equilibrium and get
  \begin{equation*}
    F_{FE} - F_{AF} \cdot \cos \ang{45} = 0 \implies F_{EF} = - \qty{15}{\kN} 
  .\end{equation*}
\end{derivation}

\finalresult{
\begin{aligned}
  F_{BC} &= \qty{18}{\kN} (\text{T}) \\
  F_{BE} &= \qty{5}{\kN} (\text{C}) \\
  F_{FE} &= \qty{15}{\kN} (\text{C})
.\end{aligned}
}

\begin{problemwithimage}{f5_36}{*5-36}
  Determine the resultant force at pins $A$, $B$, and $C$ on the two-member frame.
\end{problemwithimage}

\begin{derivation}
  We start by finding a resultant force $W$ to replace the distributed loading on $AC$ with.
  We start by calculating the length of $AC$ using trigonometry,
  \begin{equation*}
    AC = \frac{\qty{2}{\m}}{\sin \ang{60} } = \qty{2.31}{\m}
  .\end{equation*}
  Hence, the resultant of the load is
  \begin{equation*}
    W = wL = \qty{200}{\N\per\m} \cdot \qty{2.31}{\m}  = \qty{461,9}{\N}
  .\end{equation*}
  This will act perpendicular to the member meaning the force components are
  \begin{equation*}
    W_x = \qty{461.9}{\N} \cdot \cos \ang{30}  = \qty{400}{\N}, \qquad W_y = \qty{461.9}{\N} \cdot \sin (\ang{-30})  = -\qty{230.9}{\N}
  .\end{equation*}
  
  We will start by taking a global moment equilibrium about pin $A$.
  The distance from $A$ to $B$ is
  \begin{equation*}
    |AB| = \qty{2}{\m} + \cos \ang{60} \cdot \qty{2.31}{\m}  = \qty{3.155}{\meter}
  .\end{equation*}
  Hence,
  \begin{align*}
    M_{A} &= 0 \\
    \qty{800}{\N} \cdot \qty{2}{\m} + \qty{3.155}{\m} \cdot B_y - \qty{461.9}{\N} \cdot \qty{1.155}{\m}  &= 0 \\
    B_y &= \frac{\qty{461.9}{\N} \cdot \qty{1.155}{\m} - \qty{800}{\N} \cdot \qty{2}{\m} }{\qty{3.155}{\m} } \\
    B_y &= \qty{-338,04}{\N}
  .\end{align*}

  Now, we will take a force equilibrium of the whole frame in the $y$ direction as
  \begin{align*}
    F_y &= 0 \\
    A_y + B_y - W_y &= 0 \\
    A_y - \qty{338.04}{\N} - \qty{230.9}{\N} &= 0 \\
    A_y &= \qty{230.9}{\N} + \qty{338.04}{\N} \\
    A_y &= \qty{568,94}{\N}
  .\end{align*}
  
  Now, we will take a force equilibrium of the whole frame in the $x$ direction and get
  \begin{align*}
    F_x &= 0 \\
    W_x - \qty{800}{\N} + A_x + B_x &= 0 \\
    \qty{400}{\N} - \qty{800}{\N} + A_x + B_x &= 0 \\
    A_x + B_x &= \qty{400}{\N}
  .\end{align*}

  Using the method of sections, we will now cut the frame at pin $C$ and isolate the left part (member $AC$). 
  As $C$ is also a pin, the total moment here must also be zero -- Using the method of sections we further know that the total moment about $C$ from the left free body must also equal zero).
  Hence,
  \begin{align*}
    M_C &= 0 \\
    A_x \cdot \qty{2}{\m} - A_y \cdot \qty{1.155}{\m} + W \cdot \qty{1.155}{\m} &= 0 \\
    A_x &= \frac{A_y \cdot \qty{1.155}{\m} - W \cdot \qty{1.155}{\m}}{\qty{2}{\m} } \\
    A_x &= \left( \qty{568.94}{\N} - \qty{461.9}{\N} \right) \cdot \frac{\qty{1.155}{\m} }{\qty{2}{\m} } \\
    A_x &= \qty{61,78}{\N}
  .\end{align*}
  Hence,
  \begin{align*}
    A_x + B_x &= \qty{400}{\N}  \\
    B_x &= \qty{400}{\N} - \qty{61.78}{\N}  = \qty{338,22}{\N}
  .\end{align*}
  
  To get the resultant of pin $C$ we will use a force equilibrium to find the required force from the pin on the left part of the frame.
  \begin{align*}
    A_x + W_x + C_x &= 0 \\
    C_x &= -\qty{400}{\N} - \qty{61.78}{\N} = - \qty{461,78}{\N} \\
    A_y + W_y + C_y &= 0 \\
    C_y &= \qty{230.9}{\N} - \qty{568.94}{\N}   = \qty{-338,04}{\N}
  .\end{align*}
\end{derivation}

\finalresult{
\begin{aligned}
  A_x &= \qty{61,78}{\N}, & B_x &= \qty{338.22}{\N}, & C_x &= - \qty{461,78}{\N}  \\
  A_y &= \qty{568,94}{\N}, & B_y &= - \qty{338,04}{\N}, & C_y &= \qty{-338,04}{\N} 
.\end{aligned}
}
